\documentclass[10pt,a4paper]{amsart}
\usepackage[margin=19mm]{geometry}
\usepackage{amsmath,amssymb,mathtools}
\usepackage[scr=rsfs,cal=euler]{mathalfa}
\usepackage{xfrac}
\usepackage[unicode,hidelinks,bookmarks=false]{hyperref}
\newcommand{\ie}{i.e.\ }
\newcommand{\cf}{cf.\ }
\newcommand{\resp}{respectively\ }
\newcommand{\Subsection}{\subsection}
\providecommand{\nobreakdash}{\nobreak\hbox{-}}
\setlength{\emergencystretch}{2em}
\multlinegap0pt
\usepackage{bm}
\usepackage{bbm}
\usepackage{dsfont}
\usepackage{xspace}
\usepackage{cancel}
\usepackage{mathtools}
\usepackage{amsthm}
\makeatletter
\@ifundefined{thm}{\newtheorem{thm}{Thm}}{}
\@ifundefined{lem}{\newtheorem{lem}[thm]{Lemma}}{}
\makeatother
\title[The $\sigma$-inverse mean curvature flow and the generalized]{The $\sigma$-inverse mean curvature flow and the generalized Penrose conjecture}
\author{Conghan Dong}
\date{Source: arXiv:2605.26504v2 (CC BY 4.0)}
\begin{document}
\maketitle

\noindent\emph{Source and licence.} The $\sigma$-inverse mean curvature flow and the generalized Penrose conjecture. Version arXiv:2605.26504v2 (CC BY 4.0), distributed under the Creative Commons Attribution 4.0 International licence, as recorded in the arXiv licence metadata for this exact version and shown on the abstract page https://arxiv.org/abs/2605.26504v2. Full rights notice: licenses/arxiv-2605.26504-CC-BY-4.0.txt.\par\smallskip

\noindent\textit{Editorial scope.} Counted unit: the Lem whose source label is \texttt{appen-apriori-estimates} in the article (source0.tex lines 1773--1791, source label \texttt{appen-apriori-estimates}), delivered together with the article's own complete original proof of it (source0.tex lines 1793--1900, one \texttt{proof} environment of 108 lines). No mathematics is written, repaired or re-derived: statement and proof are the article's own text line for line. Required context retained uncounted: *-eq (lines 402--404); local-estimate-mean-curv (lines 433--439); appen-local-estimate-mean-curv (lines 1679--1686); epsilon-s-eq (lines 1758--1765). Every definition, display and label of those lines is retained unchanged. Omitted from this package and not redistributed: 1--401, 405--432, 440--1678, 1687--1757, 1766--1772, 1792--1792, 1901--2236. Nothing inside a retained excerpt is omitted or altered: every retained body is delivered exactly as the article's lines are written, so no omission receipt of any kind accompanies this package. Citations: the retained excerpts carry no citation command at all, so no bibliography entry of the article is transcribed and no reference is redistributed. Reference adaptation: none was needed and none was made; every reference inside the delivered unit resolves inside it, so no unresolved reference or citation is produced. Printed slips kept literal: no printed word, symbol, hypothesis or number of the counted statement or of its proof was altered. Layout and class: the article's own class is \texttt{amsart} with packages including inputenc, babel, amsmath, amssymb, amsthm, mathrsfs, esint, hyperref, xy, import, pdfpages, transparent, xcolor, comment; the wrapper is the lane's pinned \texttt{amsart} editorial wrapper at 10pt on A4, which restates the theorem-like environments the retained text needs and adds the article's own macro definitions that the retained text needs (none), with extra wrapper packages (bm, bbm, dsfont, xspace, cancel, mathtools). The delivered numbering is the unit's own. Assets: no retained span contains a figure environment, a graphics-inclusion command, a PDF-page inclusion, a table or a TikZ picture; no image file is copied into this package, the package declares no asset dependency and no image file is redistributed with it. Licence: version arXiv:2605.26504v2 of this article is distributed under the Creative Commons Attribution 4.0 International licence, as recorded in the arXiv licence metadata for this exact version and shown on the abstract page https://arxiv.org/abs/2605.26504v2. A full rights notice is delivered at licenses/arxiv-2605.26504-CC-BY-4.0.txt. Characters: every retained excerpt is pure ASCII, so the delivered engine, XeLaTeX with its default Latin Modern text font, reports no missing character.
\begin{align}\label{*-eq}
	\frac{\partial F}{\partial t}(x,t) = \frac{\nu }{H- |\nu |^2_{\sigma }}(x,t),\quad x \in N,\ 0 \leq t \leq T, \tag*{$(*)$}
\end{align} 

\begin{lem}\label{local-estimate-mean-curv}
Let $(N_t)_{0 \leq t \leq T}$ be a classical solution of \ref{*-eq} on $M^{n}$, where $N_t$ may have boundary. Then there exists a positive function $\eta (x) >0$ on $M$, depending on $g$, so that for each $x \in N_t$ and each $r< \eta(x)$, we have 
	\begin{align}
		(H - |\nu |^2_{\sigma })(x, t) \leq \max \left( (H-|\nu |^2_{\sigma })_r, \frac{C(n, r \|\sigma \|_{C^0}, r ^2 \|\nabla \sigma \|_{C^0})}{r} \right) ,
	\end{align}
	where $ (H-|\nu |^2_{\sigma })_r = \sup_{(y,s) \in P_r} (H-|\nu |^2_{\sigma })(y,s)$, and $P_r = (B_r(x) \cap N_0) \times \{0\} \cup \left( \cup _{0 \leq s \leq t}(B_r(x)\cap \partial N_s) \times \{ s\}  \right) $ is the parabolic boundary. 
\end{lem}

\begin{lem}[Lemma \ref{local-estimate-mean-curv}]\label{appen-local-estimate-mean-curv}
	Let $(N_t)_{0 \leq t \leq T}$ be a classical solution of \ref{*-eq} on $M^{n}$, where $N_t$ may have boundary. Then there exists a positive function $\eta (x) >0$ on $M$, depending on $g$, so that for each $x \in N_t$ and each $r< \eta(x)$, we have 
	\begin{align}
		(H - |\nu |^2_{\sigma })(x, t) \leq \max \left( (H-|\nu |^2_{\sigma })_r, \frac{C(n, r\|\sigma \|_{C^0}, r ^2 \|\nabla \sigma \|_{C^0})}{r} \right) ,
	\end{align}
	where $ (H-|\nu |^2_{\sigma })_r = \sup_{(y,s) \in P_r} (H-|\nu |^2_{\sigma })(y,s)$, and $P_r = (B_r(x) \cap N_0) \times \{0\} \cup \left( \cup _{0 \leq s \leq t}(B_r(x)\cap \partial N_s) \times \{ s\}  \right) $ is the parabolic boundary. 
If $(M^n, g ,\sigma  )$ is asymptotically flat, we may take $\eta(x) \geq c |x| $ for some constant $c>0$, where $x$ is the asymptotically flat coordinate.
\end{lem}

\begin{align}\label{epsilon-s-eq}
	\begin{cases}
		\mathcal{E}^{\epsilon , s}(u ^{\epsilon ,s} ):= \mathrm{div}\left( \frac{\nabla u ^{\epsilon , s}}{ \sqrt{|\nabla u^{\epsilon, s }|^2 + \epsilon ^2} } \right) - \sqrt{|\nabla u^{\epsilon, s }|^2+ \epsilon ^2} - s\cdot \frac{\sigma^{\epsilon } _{ij}\nabla ^{i} u^{\epsilon, s } \nabla ^{j} u ^{\epsilon ,s}}{|\nabla u^{\epsilon,s }|^2 + \epsilon ^2} =0 &\text{ in } \Omega _{L},\\
		u ^{\epsilon , s}=0 &\text{ on } \partial E_0,\\
		u ^{\epsilon , s}= s(L-2) & \text{ on } \partial F_{L}.
	\end{cases}
	\tag*{$(*)_{\epsilon ,s}$}
\end{align}

\begin{lem}\label{appen-apriori-estimates}
	For every $L>0$, there exists $\epsilon (L)>0$ such that for $0< \epsilon < \epsilon (L)$ and $s \in [0,1]$, a smooth solution of \ref{epsilon-s-eq} on $\bar{\Omega }_{L}$ satisfies the following a priori estimates:
	\begin{align}
		u^{\epsilon ,s} \geq - \epsilon \text{ in } \bar{\Omega }_{L},\ u^{\epsilon ,s} \geq v+ s(L-2) -L \text{ in } \bar{F}_{L}\setminus F_0,
	\end{align} 
	\begin{align}
		u^{\epsilon, s } \leq s(L-2) \text{ in } \bar{\Omega }_{L},
	\end{align}
	\begin{align}
		|\nabla u^{\epsilon ,s}| \leq H_{+}  + \epsilon \text{ on } \partial E_0,\ |\nabla u^{\epsilon ,s}| \leq C(L) \text{ on } \partial F_{L},
	\end{align}
	\begin{align}
		|\nabla u^{\epsilon ,s}(x)| \leq \max_{\partial \Omega _{L} \cap B_r(x)}|\nabla u^{\epsilon ,s}| + \epsilon + C(n, \|\sigma^{\epsilon } \|_{C^1}) \cdot r ^{-1}, \ x \in \bar{\Omega }_{L},
	\end{align}
	\begin{align}
		|u^{\epsilon ,s}|_{C^{2, \alpha }(\bar{\Omega }_{L})} \leq C(\epsilon , L, n, \|\sigma ^{\epsilon }\|_{C^1}),
	\end{align}
	for any $0<r\leq \eta (x)$. Here $H_{+}= \max(0, H_{\partial E_0})$.
\end{lem}

\begin{proof}
	Write $u = u ^{\epsilon ,s}$. 

	1. First we construct a subsolution that bridges from $E_0$ to where $v$ starts. Define $G_0 := E_0$ and $G_t := \{x: d(x, E_0) < t\} $. Select $t_L$ so that $G_{t_L} \supset F_L$. Let $\mathrm{Cut}_0$ be the cut locus of $E_0$ in $M$. On $M \setminus (E_0 \cup \mathrm{Cut}_0)$, $d(\cdot , E_0)$ is smooth, and each point is connected to $E_0$ by a unique minimizing geodesic $\gamma $, and $\partial G_t$ foliates a neighborhood of $\gamma $. We have
	\begin{align*}
		\frac{\partial H}{\partial t} &= - |\mathrm{II}|^2 - \mathrm{Ric}(\nu ,\nu ) \leq C_1(L)\quad  \text{ on } \partial G_t \setminus \mathrm{Cut}_0, \ 0 \leq t \leq t_L,
	\end{align*} 
	yielding
	\begin{align*}
		H_{\partial G_t} \leq \max_{\partial E_0} H_{+} + C_1 t \leq C_2(L) \quad  \text{ on } \partial G_t \setminus \mathrm{Cut}_0, \ 0 \leq t \leq t_L,
	\end{align*} 
	where $H_{+} = \max (0, H)$.  Consider the prospective subsolution
	\begin{align*}
		v_1(x) := f(x) = f(d(x, G_0) ),\quad x \in \bar{G}_{t_L}\setminus E_0,
	\end{align*} 
	with $f'<0$. Then $\nabla v_1 = f' \nu , \nabla ^2 v_1 = f'' \nu \otimes \nu + f' \mathrm{II}$, so
	\begin{align*}
		\Delta v_1 - \nabla ^2 v_1(\nu ,\nu ) = f' H_{\partial G_t} \geq - C_2|f'| .
	\end{align*} 
	Hence
	\begin{align*}
		&\sqrt{|f'|^2+ \epsilon ^2} \mathcal{E}^{\epsilon ,s}(v_1) \\
		&= \Delta v_1 -  \left( |f'|^2 + \epsilon ^2 \right) ^{-1} |f'|^2 f'' - |f'|^2- \epsilon ^2 - \left( |f'|^2+ \epsilon ^2 \right) ^{-\frac{1}{2}} s|f'|^2 \sigma ^{\epsilon } (\nu ,\nu ) \\
		&\geq -C_2 |f'| + \frac{\epsilon ^2}{|f'|^2+ \epsilon ^2} f''- |f'|^2- \epsilon ^2 - \|\sigma ^{\epsilon } \|_{C^0}\cdot |f'|.
	\end{align*} 
	Set
	\begin{align*}
		f(t) := \frac{\epsilon }{A}(-1 + e ^{-A t}),\quad 0 \leq t \leq t_L	.
	\end{align*} 
	Then we have $\epsilon ^2 \leq |f'| = \epsilon e^{-At} \leq \epsilon $ and $f''  = A |f'|$, provided that we impose $\epsilon \leq \epsilon (A,L):= e^{-A t_L}$. Choosing $A:= 2(C_2+\|\sigma^{\epsilon } \|_{C^0}+2)$, we have
	\begin{align*}
		(|f'|^2+ \epsilon ^2) \left( C_2|f'| + |f'|^2 + \epsilon ^2 + \|\sigma^{\epsilon } \|_{C^0} |f'| \right) \leq 2 \epsilon ^2 \left( C_2 + \epsilon + 1+ \|\sigma^{\epsilon } \|_{C^0} \right)  |f'| \leq \epsilon ^2 f''. 
	\end{align*} 
	This shows that the function
	$$v_1(x) := \frac{\epsilon }{2(C_2+\|\sigma ^{\epsilon } \|_{C^0} + 2)} \left( -1 + e^{- 2(C_2+\|\sigma ^{\epsilon } \|_{C^0}+2) d(x, \partial G_0)} \right) ,$$ 
	is a smooth subsolution for $\mathcal{E}^{\epsilon ,s}$ on $G_{t_L} \setminus (E_0 \cup \mathrm{Cut}_0)$ for sufficiently small $\epsilon $.

	It can be shown that $v_1$ is a viscosity subsolution of $\mathcal{E}^{\epsilon ,s}$ on all of $G_{t_L}\setminus \bar{E}_0$. Since $u \geq v_1$ on the boundary, it follows by the maximum principle for viscosity solution that 
	\begin{align}\label{v1-lower}
		u \geq v_1 \geq -\epsilon \text{ in } \bar{\Omega }_{L},\ \frac{\partial u}{\partial \nu } \geq - \epsilon \text{ on }\partial E_0.
	\end{align} 

2. Next, consider the function
\begin{align*}
	v_2:= \frac{L-1}{L} v + s(L-2) - (L-1).
\end{align*} 
Then $\mathcal{E}^{0,s}(v_2) = \mathcal{E}^{0,s}(v) + \frac{1}{L}|\nabla v| >0$ on $\bar{F}_{L} \setminus F_0$, where we used the assumption that $\mathcal{E}^{0,s}(v) \geq \mathcal{E}^{0,1}(v) \geq 0$ outside $F_0$. Since the domain is compact, for all sufficiently small $\epsilon $ we obtain $\mathcal{E}^{\epsilon ,s}(v_2) >0$. Note that
\begin{align*}
	u \geq - \epsilon \geq v_2 \text{ on } \partial F_0,\quad u= s(L-2) = v_2 \text{ on } \partial F_L.
\end{align*} 
By the maximum principle, 
\begin{align}\label{v2-lower}
	u \geq v_2 \geq v+ s(L-2) - L \text{ in } \bar{F}_{L}\setminus F_0,\quad \frac{\partial u}{\partial \nu } \geq - C(L) \text{ on } \partial F_L.
\end{align}

Since $\mathcal{E}^{\epsilon ,s}(s(L-2) ) = - \epsilon \leq 0$, we obtain
\begin{align}
	u \leq s(L-2) \text{ in } \bar{\Omega }_{L},\quad \frac{\partial u}{\partial \nu } \leq 0 \text{ on } \partial F_{L}.
\end{align}

3. Next we construct a supersolution along $\partial E_0$. Choose a smooth function $v_3$ which vanishes on $\partial E_0$ such that 
\begin{align*}
	H_{+} < \frac{\partial v_3}{\partial \nu } \leq H_+  + \epsilon \quad \text{ along } \partial E_0.
\end{align*}
Note that along $\partial E_0$, 
\begin{align*}
	\mathcal{E}^{0, s}(v_3) = H_{\partial E_0} - |\nabla v_3| - s \cdot \sigma ^{\epsilon } (\nu ,\nu ) \leq H_{+}   - |\nabla v_3| <0, 
\end{align*}
This implies that for sufficiently small $\delta >0$, $|\nabla v_3|>0$ and $\mathcal{E}^{0,s}(v_3)<0$ in the neighborhood $U:= \{ 0 \leq v_3 \leq \delta \} $. Now define the sped-up function
 \begin{align*}
	v_4 := \frac{v_3}{1- v_3 / \delta },\quad x \in U.
\end{align*} 
Then $v_4 \to \infty$ on $\partial U \setminus \partial E_0$, and $\mathcal{E}^{0,s}(v_4) = \mathcal{E}^{0,s}(v_3) + |\nabla v_3|\left( 1- (1- v_3 / \delta )^{-2} \right) <0 $. For sufficiently small $\epsilon $, depending on $L$, $\mathcal{E}^{\epsilon ,s}(v_4) <0 $ on the set $V:= \{0 \leq v_4 \leq L\} $. Since $u \leq L-2$ in $\bar{\Omega }_{L}$, we have $u \leq v_4$ on $\partial V$. By the maximum principle, we obtain $u \leq v_4$ on $V$, and therefore
\begin{align}\label{v4-upper}
	\frac{\partial u}{\partial \nu } \leq \frac{\partial v_4}{\partial \nu }= \frac{\partial v_3}{\partial \nu } \leq H_+ + \epsilon \quad \text{ on } \partial E_0.
\end{align} 

4. Let $\tilde{N}_t ^{\epsilon ,s}$ denote the level-set $\{U = t\} \subset \Omega _{L} \times \mathbb{R} $ of the function $U(x,z)=U_{\epsilon ,s}(x, z) := u^{\epsilon ,s}(x) - \epsilon z, -\infty< t < \infty$, where $\Omega _{L} \times \mathbb{R}$ is equipped with the product metric $\tilde{g}= g + d z^2$. Equation \ref{epsilon-s-eq} is equivalent to
\begin{align*}
	\mathrm{div}_{\tilde{g}}\left( \frac{\tilde{\nabla } U}{|\tilde{\nabla } U|} \right)  = |\tilde{\nabla }U| + s\cdot  \frac{\sigma _{ij} ^{\epsilon }\tilde{\nabla}^i U \tilde{\nabla} ^j U}{|\tilde{\nabla} U|^2 },
\end{align*}
where we abused the notation $\sigma ^{\epsilon } (x,z) = \sigma ^{\epsilon } (x)$ to denote the constant extension of $\sigma ^{\epsilon } $ along $z$-direction. By the level-set description, $\tilde{N}_t ^{\epsilon ,s}$ is a smooth solution of the $(s \sigma^{\epsilon }) $-IMCF:
\begin{align*}
	\frac{\partial \tilde{F}}{\partial t} = \frac{\tilde{\nu }}{\tilde{H}- |\tilde{\nu }|^2_{s \sigma ^{\epsilon } }},
\end{align*} 
with $\tilde{\nu } = \frac{\tilde{\nabla }U}{|\tilde{\nabla }U|}$ and $\tilde{H} = \mathrm{div}_{\tilde{g}}\left( \frac{\tilde{\nabla } U}{|\tilde{\nabla } U|} \right) $.  
Applying Lemma \ref{appen-local-estimate-mean-curv} to $\tilde{N}^{\epsilon ,s}_t$, we obtain the estimate
\begin{align*}
	 \tilde{H} - |\tilde{\nu }|^2_{s\sigma ^{\epsilon } } \leq \max \left( (\tilde{H}-|\tilde{\nu }|^2_{s\sigma ^{\epsilon } })_{\tilde{P}_r}, \frac{C(n, \|\sigma ^{\epsilon } \|_{C^1})}{r} \right) .
\end{align*}
Notice that $\tilde{H} - |\tilde{\nu }|^2_{s\sigma ^{\epsilon } } = |\tilde{\nabla }U| = \sqrt{|\nabla u |^2+ \epsilon ^2}  $ is independent of $z$. We obtain
\begin{align*}
	\sqrt{|\nabla u |^2+ \epsilon ^2}(x) &\leq \sup_{t} \max_{\partial \tilde{N}^{\epsilon ,s}_{t} \cap \tilde{B}_r(x,z)} \sqrt{|\nabla u |^2+ \epsilon ^2} + \frac{C(n, \|\sigma ^{\epsilon } \|_{C^1})}{r} \\
	&\leq \max_{\partial \Omega _{L} \cap B_r(x)} |\nabla u| + \epsilon + \frac{C}{r},
\end{align*} 
for $x \in M \setminus E_0$ and any $0<r \leq \eta _M(x)$. 

Thus we have the Lipschitz estimate
\begin{align*}
	|u|_{C^{0,1}(\bar{\Omega }_{L})} \leq C(L, n, \|\sigma ^{\epsilon } \|_{C^1}).
\end{align*} 
Following the Nash-Moser-De Giorgi estimates, for some $\alpha =\alpha (\Omega _{L})$, we obtain
\begin{align*}
	\|u\|_{C^{1, \alpha }(\bar{\Omega }_{L})} \leq C(\epsilon ,L,n, \|\sigma ^{\epsilon } \|_{C^1}).
\end{align*} 
The Schauder estimates complete the proof.

\end{proof}

\end{document}
